\begin{abstract}
Many LLM and agent papers introduce workflows that could help ordinary users and
agent builders, but those workflows are usually described for human reading
rather than direct reuse. Summaries compress the idea, but often discard the
procedural details, validation checks, assumptions, and failure branches needed
to apply the method. We introduce PaperToSkill, a paper-to-skill conversion
workflow that turns source-anchored paper notes into compact, human-editable
agent skills. Each skill records the paper's contribution, use conditions,
workflow, validation checks, failure cases, transfer notes, and source anchors.
In a four-paper benchmark over AI Scientist-v2, Reflexion, AIDE, and
Toolformer, PaperToSkill generates skills that score 20/20 on deterministic
structural rubrics, preserve more deterministic operational coverage than
generic-summary and abstract-only baselines, remain under a 1200-word
compactness budget, and substantially reduce deterministic input-token proxy
size relative to full extracted papers. These results support PaperToSkill as a
reproducible conversion layer from curated paper notes to portable skills. All
four live-transfer saved-response sets are collected and scored under a
deterministic output-contract evaluator; these saved response files are not
human semantic fidelity or downstream execution-outcome evidence. A first
single-run real-reuse stress test over eight AIDE, SWE-agent, Reflexion, and
SnapATAC2 paper-task rows is complete, but the outcome is mixed rather than
confirmatory: AIDE-T2 and one SWE task favor PaperToSkill, AIDE-T1 and
Reflexion are solved by both conditions, and SWE-T1 plus SnapATAC2 expose
patch-application and artifact-completion failures. We therefore treat these
rows as downstream boundary evidence, not aggregate effectiveness.
\end{abstract}
